\begin{center}
\begin{tikzpicture}[scale=0.75,every node/.style={circle,draw=black!50,minimum size=0.48cm,inner sep=0pt},line width=0.6pt]
\node[fill=newsblue!65] (a) at (0,0) {};
\node[fill=orange!65] (b) at (1.4,0) {};
\node[fill=newsblue!65] (c) at (2.8,0) {};
\node[fill=orange!65] (d) at (4.2,0) {};
\draw (a)--(b)--(c)--(d);
\node[draw=none,font=\scriptsize] at (2.1,-0.7) {$P_4$：路径};
\node[fill=newsblue!65] (e) at (7.1,0) {};
\node[fill=orange!65] (f) at (7.1,1.2) {};
\node[fill=orange!65] (g) at (5.8,-0.75) {};
\node[fill=orange!65] (h) at (8.4,-0.75) {};
\draw (e)--(f) (e)--(g) (e)--(h);
\node[draw=none,font=\scriptsize] at (7.1,-1.35) {$K_{1,3}$：星形图};
\end{tikzpicture}
\par\footnotesize\color{newsgray}图3　编辑部原创二色正常染色示例。两图都能用两色，但连接结构不同；本图不是论文计算实验的原始图。\end{center}